\section{Introduction}

Periodic points of polynomial automorphisms carry two different kinds of
structure.  Their real dynamics may admit a symbolic coding and uniform
hyperbolicity, while their complex coordinates form finite algebraic
schemes.  Ordinarily these descriptions need not exhaust one another: a
real horseshoe can coexist with additional complex periodic points, and an
algebraic intersection can contain multiplicity invisible to symbolic
dynamics.

This paper identifies a parameter at which the two ledgers coincide exactly.
Consider
\begin{equation}\label{eq:H6}
H_6(q,p)=(1-6q^2-p,q).
\end{equation}
The map is a degree-two polynomial automorphism with Jacobian determinant
one.  Our main result is the following.

\begin{theorem}[Algebraic exhaustion]\label{thm:exhaustion-intro}
For every integer $n\ge1$, the complex fixed-point scheme of $H_6^n$ is
reduced and consists of exactly $2^n$ real uniformly hyperbolic points.
Equivalently,
\[
\Fix_{\mathbb C}(H_6^n)=\Fix_{\mathbb R}(H_6^n),
\qquad \#\Fix(H_6^n)=2^n.
\]
\end{theorem}

The proof is not a large finite computation.  It is a short bridge among
three source-native theorems.  First, a linear conjugacy identifies
\eqref{eq:H6} with the standard area-preserving H\'enon map at $a=6$.
Arai's rigorous hyperbolic plateau contains the entire parameter path from
$6$ to $10$ \cite{Arai2007}.  At the latter endpoint the
Devaney--Nitecki large-parameter theorem gives a full two-shift
\cite{DevaneyNitecki1979}.  Hyperbolic structural stability transports this
symbolic system to $a=6$, yielding $2^n$ distinct real fixed points.  Second,
Friedland--Milnor's intersection theorem says that the total complex
algebraic multiplicity is also $2^n$ \cite{FriedlandMilnor1989}.  The real
points therefore consume the entire complex count.

The motivation comes from a reflection slice developed in two predecessor
projects.  If $n=2m+1$ is odd, an exact mixed-axis closure has the form
\begin{equation}\label{eq:Fn-intro}
F_n(X)=q_{m+1}(X)-q_m(X),\qquad \deg F_n=2^{m+1},
\end{equation}
where the $q_j$ are generated by the H\'enon recurrence on one reversor-fixed
line.  Formal M\"obius division predicts a primitive degree of entropy
$\tfrac12\log2$.  Previous work established divisibility and transversality
only on a smaller certified survivor.  Theorem~\ref{thm:exhaustion-intro}
removes the ambient gap.

\begin{corollary}[Totally real primitive reflection divisors]
For every odd $n$, $F_n$ is totally real and squarefree.  Its recursive
M\"obius quotient $\Psi_n$ is a reduced effective divisor supported exactly
on least-period-$n$ points.  In particular,
\[
\mathbb Q[X]/(\Psi_n)
\]
is a finite \emph{\'etale} totally real algebra.
\end{corollary}

The result is deliberately not phrased as an arithmetic spectral theorem.
The total reality and exact entropy provide a stronger all-period object for
future height and Galois-pressure questions, but they produce neither
rational-prime labels nor the prime-power amplitudes required by a Riemann
trace formula.  The claim boundary is part of the theorem package.

Section~\ref{sec:sources} locks parameter conventions and source scopes.
Section~\ref{sec:horseshoe} transports the full horseshoe, and
Section~\ref{sec:exhaustion} proves complex algebraic exhaustion.
Sections~\ref{sec:reflection}--\ref{sec:entropy} derive reflection
effectivity and entropy.  Section~\ref{sec:finite} reports the exact audit,
and Section~\ref{sec:boundary} records the remaining arithmetic boundary.
